%! AP Statistics — Unit 9, Lesson 9.2 — Homework
\documentclass[10pt]{article}
\usepackage{apstats-article}
\usepackage{apstats-boxes}

\begin{document}

\pageheader{Unit 9, Lesson 9.2}{Homework}
\namedateperiod

\vspace{0.1in}

\textbf{Problem 1: Sleep and Reaction Time}

\vspace{0.05in}
A random sample of 22 college students was used to study the relationship between hours
of sleep per night ($x$) and reaction time in milliseconds ($y$). Computer output is
shown below.

\vspace{0.1in}
\begin{center}
\renewcommand{\arraystretch}{1.3}
\begin{tabular}{lrrrr}
    \textbf{Predictor} & \textbf{Coef} & \textbf{SE Coef} & \textbf{T} & \textbf{P} \\
    \hline
    Constant & 412.3 & 28.4 & 14.52 & 0.000 \\
    Sleep    & $-18.6$ & 5.3 & $-3.51$ & 0.002 \\
    \hline
    \multicolumn{5}{l}{$s = 24.1$ \quad $n = 22$}
\end{tabular}
\end{center}

\vspace{0.1in}
\begin{enumerate}[label=\alph*.]

    \item Identify $b$ and $SE_b$ from the output.
          \writelines{1}

    \item For this problem, assume: the data come from a random sample; the residual plot
          shows random scatter with roughly constant spread; and $n = 22$ is large enough
          for normality (stated). Briefly note how each LINER condition is addressed.
          \writelines{5}

    \item Construct a 95\% CI for $\beta$ using $t^* = 2.086$ ($df = 20$). Show all work.
          \writelines{3}

    \item Interpret the CI in context.
          \writelines{2}

    \item Does the CI provide evidence of a negative linear relationship between hours of
          sleep and reaction time? Explain.
          \writelines{2}

\end{enumerate}

\vspace{0.2in}

\textbf{Problem 2: Temperature and Heating Fuel Cost}

\vspace{0.05in}
A researcher studies 18 cities, recording average daily high temperature in January
($x$, ${}^\circ$F) and monthly heating fuel cost ($y$, \$/month).
Output gives $b = -4.8$, $SE_b = 1.1$, and $n = 18$.
Use $t^* = 2.120$ ($df = 16$). Assume all LINER conditions are met.

\vspace{0.1in}
\begin{enumerate}[label=\alph*.]

    \item Construct a 95\% CI for the true slope $\beta$. Show all work, then interpret
          the interval in context.
          \writelines{4}

    \item A utility company claims $\beta = -5$ (each 1${}^\circ$F warmer costs \$5 less
          per month). Is this claim consistent with your CI? Explain.
          \writelines{2}

    \item If the study had used 72 cities instead of 18, would you expect the CI to be
          narrower or wider? Explain using the formula for $SE_b$.
          \writelines{2}

\end{enumerate}

\end{document}
